% BOOK-FILE-ID: f-de418d1ee182
% 付録：行列の基本

\section{行列の基本}\label{節：行列の基本}

\sectionref*{節：連立型}では,二つの\bookterm{sequence}{数列}から和と差を作り,漸化式を簡単にした.
例えば
\[
  x_{n+1}=2x_n+y_n,\qquad y_{n+1}=x_n+2y_n
\]
なら,\,$(x_1,y_1)=(1,0)$から$(x_2,y_2)=(2,1)$,次に$(x_3,y_3)=(5,4)$と進む.
毎回の計算には2つの値が両方必要なので,いっそ2つの値を1組として記録し,その組をまとめて計算してしまおう.
そのとき,漸化式の係数をまとめて書くための道具が\bookterm{matrix}{行列}である.
まずは行列の掛け方を学び,そこから連立型で行った「等比型への分解」が何をしていたのかを見ていく.

\subsection{行列とベクトル}

\bookterm{matrix}{行列}とは,数を長方形に並べ,全体を括弧で囲んだものである.
たとえば
\[
  A=\begin{pmatrix}2&1\\1&2\end{pmatrix}
\]
は二行二列の\bookterm{matrix}{行列}である.横の並びを\textbf{行},縦の並びを\textbf{列},各数を\textbf{成分}という.
行は上から,列は左から数え,この例の第一行・第二列の成分は$1$である.
行列の定数倍は各成分をその定数倍する.

また,二つの数を縦に並べた
\[
  \boldsymbol{v}=\begin{pmatrix}x\\y\end{pmatrix}
\]
を\textbf{\bookterm{column-vector}{列ベクトル}}と呼ぶ.
太字の$\boldsymbol{v}$は二つの値をまとめた記号で,平面上の点$(x,y)$への矢印とも読める.
0でない二つのベクトルが直交するとは,矢印が垂直に交わることである.
成分が$(u_1,u_2)$と$(v_1,v_2)$なら,直交する条件は$u_1v_1+u_2v_2=0$である.

\bookterm{column-vector}{列ベクトル}の足し算と定数倍は,成分ごとに行う.
\[
  \begin{pmatrix}x\\y\end{pmatrix}
  +\begin{pmatrix}u\\v\end{pmatrix}
  =\begin{pmatrix}x+u\\y+v\end{pmatrix},
  \qquad
  c\begin{pmatrix}x\\y\end{pmatrix}
  =\begin{pmatrix}cx\\cy\end{pmatrix}.
\]
たとえば$\begin{pmatrix}3\\1\end{pmatrix}+\begin{pmatrix}1\\-1\end{pmatrix}=\begin{pmatrix}4\\0\end{pmatrix}$である.
等号も成分ごとの等しさを表す.
一般の二成分のベクトルも
\[
  \begin{pmatrix}x\\y\end{pmatrix}
  =\frac{x+y}{2}\begin{pmatrix}1\\1\end{pmatrix}
  +\frac{x-y}{2}\begin{pmatrix}1\\-1\end{pmatrix}
\]
と分けられる.本編の和と差は,この二つの係数を求める操作に対応する.
\bookproblemref{practice-basic}{7}{標準問7}でも使うこの分解が,後で固有値に結び付く.

\subsubsection{行列を掛ける}

先ほどの$A$を\bookterm{column-vector}{列ベクトル}に掛けると
\[
  \begin{pmatrix}2&1\\1&2\end{pmatrix}
  \begin{pmatrix}x\\y\end{pmatrix}
  =
  \begin{pmatrix}2x+y\\x+2y\end{pmatrix}.
\]
となる.各行の数を$x,y$にそれぞれ掛けて足す.
たとえば
\[
  A\begin{pmatrix}1\\0\end{pmatrix}
  =\begin{pmatrix}2\\1\end{pmatrix},
  \qquad
  A\begin{pmatrix}2\\1\end{pmatrix}
  =\begin{pmatrix}5\\4\end{pmatrix}.
\]
$A$は$x,y$を$2x+y,x+2y$へ移す規則を表す.
\BookFigMatrixMap

一般には
\[
  \begin{pmatrix}p&q\\r&s\end{pmatrix}
  \begin{pmatrix}x\\y\end{pmatrix}
  =\begin{pmatrix}px+qy\\rx+sy\end{pmatrix}
\]
と定める.連立式$x'=px+qy,\ y'=rx+sy$を積一つにまとめたもので,\,$x',y'$は$A$を掛けた後の値を表す.
このように行列を掛けてベクトルを移す操作を一次変換という.

また,成分ごとに展開すれば
\[
  A(c\boldsymbol{u}+d\boldsymbol{v})
  =cA\boldsymbol{u}+dA\boldsymbol{v}
\]
が成り立つ.
つまり,足してから$A$を掛けても,別々に$A$を掛けてから足しても同じになる.後で初期値を分ける際に使う.

\subsection{行列の積}

$A$を2回掛ける操作をまとめた\bookterm{matrix}{行列}を$A^2$と書く.文字$x,y$で計算すると
\[
  \begin{aligned}
    A\left(A\begin{pmatrix}x\\y\end{pmatrix}\right)
    &=A\begin{pmatrix}2x+y\\x+2y\end{pmatrix}\\
    &=\begin{pmatrix}5x+4y\\4x+5y\end{pmatrix}.
  \end{aligned}
\]
したがって$A^2=\begin{pmatrix}5&4\\4&5\end{pmatrix}$である.各成分の二乗とは異なる.

異なる\bookterm{matrix}{行列}$A,B$でも,まず$B$,次に$A$を掛ける操作を$AB$と書く.
$B$の各列に$A$を掛けた結果を並べると,積$AB$が得られる.
具体的には
\[
  \begin{aligned}
    &\begin{pmatrix}p&q\\r&s\end{pmatrix}
    \begin{pmatrix}u&v\\w&z\end{pmatrix}\\
    &\qquad=
    \begin{pmatrix}
      pu+qw&pv+qz\\
      ru+sw&rv+sz
    \end{pmatrix}.
  \end{aligned}
\]
右の\bookterm{matrix}{行列}から先に作用するため,一般には$AB$と$BA$は一致しない.

何も変えない変換を表す行列は
\[
  I=\begin{pmatrix}1&0\\0&1\end{pmatrix},
  \qquad
  I\begin{pmatrix}x\\y\end{pmatrix}
  =\begin{pmatrix}x\\y\end{pmatrix}
\]
で表される.
$I$を\textbf{\bookterm{identity}{単位行列}}という.
数の累乗で$a^0=1$とするのに対応して,\bookterm{matrix}{行列}では$A^0=I$と定める.

\subsection{回転と複素数の掛け算}\label{節：回転と複素数の掛け算}

同じ\bookterm{matrix}{行列}を繰り返し掛けると何が起こるのか,図で見やすい例を1つ用意しておこう.
原点を中心に,点を角$\theta$だけ回す操作である.

長さ1の点$(\cos\varphi,\sin\varphi)$を角$\theta$だけ回すと,\,$(\cos(\varphi+\theta),\sin(\varphi+\theta))$に移る.
\bookterm{addition}{加法定理}で展開すると
\[
  \begin{aligned}
    \cos(\varphi+\theta)&=\cos\theta\cos\varphi-\sin\theta\sin\varphi,\\
    \sin(\varphi+\theta)&=\sin\theta\cos\varphi+\cos\theta\sin\varphi
  \end{aligned}
\]
である.\,$\cos\varphi,\sin\varphi$に掛かる係数を並べると
\[
  R(\theta)=\begin{pmatrix}\cos\theta&-\sin\theta\\\sin\theta&\cos\theta\end{pmatrix}
\]
となり,\,$R(\theta)\begin{pmatrix}\cos\varphi\\\sin\varphi\end{pmatrix}=\begin{pmatrix}\cos(\varphi+\theta)\\\sin(\varphi+\theta)\end{pmatrix}$と書ける.
長さ$r$の点$r(\cos\varphi,\sin\varphi)$も,両辺を$r$倍すれば同じように回る.
この$R(\theta)$を\textbf{回転行列}という.

回転を2回続ければ角$2\theta$の回転になるので,\,$R(\theta)^2=R(2\theta)$である.
繰り返せば
\[
  R(\theta)^n=R(n\theta)
\]
となる.同じ行列を繰り返し掛けることが,一定の角を何度も足すことになるのだ.

この行列は,複素数の掛け算にも現れる.
複素数$a+bi$を$x+yi$に掛けると
\[
  (a+bi)(x+yi)=(ax-by)+(bx+ay)i
\]
となる.実部と虚部を縦に並べたベクトルで見ると
\[
  \begin{pmatrix}ax-by\\bx+ay\end{pmatrix}
  =\begin{pmatrix}a&-b\\b&a\end{pmatrix}\begin{pmatrix}x\\y\end{pmatrix}
\]
である.ここで$a+bi=r(\cos\theta+i\sin\theta)$と極形式で書けば,
\[
  \begin{pmatrix}a&-b\\b&a\end{pmatrix}=rR(\theta)
\]
となる.つまり,複素数を掛けることは,\textbf{角$\theta$の回転と$r$倍の拡大・縮小}を表す行列を掛けることなのである.
\sectionref*{節：特性方程式の解が虚数の場合}で複素数の冪が回転として現れたのは,このためだったのだ.

\subsection{固有ベクトルと固有値}

毎回\bookterm{matrix}{行列}の積を計算するのは手間がかかる.
そこで,掛けても定数倍になるだけの\bookterm{column-vector}{列ベクトル}を探そう.
先ほどの$A$について
\[
  \begin{aligned}
    A\begin{pmatrix}1\\1\end{pmatrix}
    &=\begin{pmatrix}3\\3\end{pmatrix}
    =3\begin{pmatrix}1\\1\end{pmatrix},\\
    A\begin{pmatrix}1\\-1\end{pmatrix}
    &=\begin{pmatrix}1\\-1\end{pmatrix}
  \end{aligned}
\]
である.
一方は毎回$3$倍され,他方は変わらないので
\[
  A^k\begin{pmatrix}1\\1\end{pmatrix}
  =3^k\begin{pmatrix}1\\1\end{pmatrix},
  \qquad
  A^k\begin{pmatrix}1\\-1\end{pmatrix}
  =\begin{pmatrix}1\\-1\end{pmatrix}
\]
となる.

一般に,0でない\bookterm{column-vector}{列ベクトル}$\boldsymbol{u}$が
\[
  A\boldsymbol{u}=\lambda\boldsymbol{u}
\]
を満たすとき,\,$\boldsymbol{u}$を$A$の\textbf{\bookterm{eigenvector}{固有ベクトル}},\,$\lambda$を対応する\textbf{\bookterm{eigenvalue}{固有値}}という.
\bookterm{eigenvalue}{固有値}は,そのベクトルに対して$A$を1回掛けると何倍になるかを表す数である.
固有ベクトルを0でない数で定数倍したものも,同じ固有値の固有ベクトルである.原点を通る直線(方向)だけが決まり,長さや向きは自由に選べる.
零ベクトルはどんな$\lambda$でも式を満たすため除く.負の\bookterm{eigenvalue}{固有値}では向きが反転し,\,$0$なら1回掛けただけで零ベクトルになる.

\subsubsection{固有ベクトルに分ける}

\bookterm{eigenvector}{固有ベクトル}が見つかると,ほかのベクトルへの作用も計算しやすくなる.
たとえば
\[
  \begin{pmatrix}1\\0\end{pmatrix}
  =\frac12\begin{pmatrix}1\\1\end{pmatrix}
  +\frac12\begin{pmatrix}1\\-1\end{pmatrix}
\]
である.
\booknavref{解釈：三項間と解全体}{隣接3項型}で一般項を2つの解の和に分けたように,ここでは$(1,0)$というベクトルを,二つの\bookterm{column-vector}{列ベクトル}の\bookterm{combination}{線形結合}で表している.
2つの部分に別々に$A$を$k$回掛ければ
\[
  \begin{aligned}
    A^k\begin{pmatrix}1\\0\end{pmatrix}
    &=\frac{3^k}{2}\begin{pmatrix}1\\1\end{pmatrix}
    +\frac12\begin{pmatrix}1\\-1\end{pmatrix}\\
    &=\begin{pmatrix}\frac{3^k+1}{2}\\\frac{3^k-1}{2}\end{pmatrix}.
  \end{aligned}
\]
本編の等比型への帰着と同じく,一定の倍数で変わる部分に分ければ,大きな累乗も計算できる.
図では,定数倍される二つのベクトルを$\boldsymbol v,\boldsymbol w$,
一般の方向の例を$\boldsymbol u=(1,0)$として描く.
\BookFigEigen

一般に,二つの異なる\bookterm{eigenvalue}{固有値}$\lambda_1,\lambda_2$に対応する\bookterm{eigenvector}{固有ベクトル}$\boldsymbol{u}_1,\boldsymbol{u}_2$があれば,どんな二成分のベクトルも
\[
  \boldsymbol{v}=c\boldsymbol{u}_1+d\boldsymbol{u}_2
\]
と表せる.
異なる\bookterm{eigenvalue}{固有値}に対応する二つの\bookterm{eigenvector}{固有ベクトル}は,互いの定数倍にならないからである.
係数$c,d$は,上下の成分を比較した連立方程式で求められる.
その後は
\[
  A^k\boldsymbol{v}
  =c\lambda_1^k\boldsymbol{u}_1+d\lambda_2^k\boldsymbol{u}_2
\]
で計算できる.
$k=0$のときは元の分解式を読み,\,$k\ge1$で\bookterm{eigenvalue}{固有値}が$0$なら対応する部分は0になる.

\subsection{行ベクトルと左固有ベクトル}

数を横に並べた$L=(\alpha,\beta)$を\textbf{\bookterm{row-vector}{行ベクトル}}という.
\bookterm{column-vector}{列ベクトル}$\boldsymbol{\ell}$を横に並べ直す操作を\textbf{\bookterm{transpose}{転置}}といい,\,$\boldsymbol{\ell}^{\mathsf T}$と書く.
たとえば$\boldsymbol{\ell}=\begin{pmatrix}\alpha\\\beta\end{pmatrix}$なら,\,$\boldsymbol{\ell}^{\mathsf T}=(\alpha,\beta)$である.
行列の転置も,行と列を入れ替える操作であり,\,$A^{\mathsf T}$と書く.
\[
  \begin{pmatrix}p&q\\r&s\end{pmatrix}^{\mathsf T}
  =\begin{pmatrix}p&r\\q&s\end{pmatrix}.
\]
\bookterm{row-vector}{行ベクトル}と\bookterm{column-vector}{列ベクトル}の積は,対応する成分を掛けて足す.
\[
  (\alpha,\beta)\begin{pmatrix}x\\y\end{pmatrix}
  =\alpha x+\beta y.
\]
\bookterm{row-vector}{行ベクトル}を\bookterm{matrix}{行列}に掛けるときは,\bookterm{matrix}{行列}の各列との積を横に並べる.
たとえば,先ほどの$A$について
\[
  \begin{aligned}
    (1,\,1)A&=(3,\,3)=3(1,\,1),\\
    (1,-1)A&=(1,-1)
  \end{aligned}
\]
である.
一般に,0でない\bookterm{row-vector}{行ベクトル}$L$が$LA=\lambda L$を満たすとき,\,$L$を\textbf{\bookterm{left-eigen}{左固有ベクトル}}という.
この場合,\bookterm{column-vector}{列ベクトル}$\boldsymbol{v}$から作った数$L\boldsymbol{v}$は,\,$A$を掛けた後には
\[
  L(A\boldsymbol{v})=(LA)\boldsymbol{v}
  =\lambda L\boldsymbol{v}
\]
となる.
たとえば$L=(1,\,1)$なら,\,$x+y$という和が$A$を掛けるたびに$3$倍される.
\bookterm{eigenvector}{固有ベクトル}は定数倍になる\bookterm{column-vector}{列ベクトル}を表し,\bookterm{left-eigen}{左固有ベクトル}は定数倍になる量を取り出す係数を表す.

\subsection{行列式と固有値}

\bookterm{eigenvalue}{固有値}の条件をまとめるため,\bookterm{matrix}{行列}式を導入する.
二行二列の場合には
\[
  \det\begin{pmatrix}a&b\\c&d\end{pmatrix}=ad-bc
\]
と定める.
\bookterm{matrix}{行列}そのものではなく,四つの成分から計算した一つの数である.

二つの式$ax+by=0,\ cx+dy=0$について,\,$ad-bc\ne0$なら解は$x=y=0$だけになる.
たとえば,第一の式を$d$倍したものから第二の式を$b$倍したものを引くと,\,$(ad-bc)x=0$となる.
同様に$(ad-bc)y=0$も得られる.
逆に$ad-bc=0$なら,0でない解が存在する.
実際,\,$(a,b)\ne(0,\,0)$なら$(x,y)=(b,-a)$が両方の式を満たす.
$(a,b)=(0,\,0)$の場合も,残る一つの式から0でない解を選べる.

\bookterm{eigenvalue}{固有値}の式$A\boldsymbol{u}=\lambda\boldsymbol{u}$は
\[
  (\lambda I-A)\boldsymbol{u}=\boldsymbol{0}
\]
と書き直せる.
\bookterm{matrix}{行列}の引き算は対応する成分ごとに行うので,\,$A=\begin{pmatrix}p&q\\r&s\end{pmatrix}$なら
\[
  \lambda I-A
  =\begin{pmatrix}\lambda-p&-q\\-r&\lambda-s\end{pmatrix}.
\]
0でない解が存在する条件は
\[
  \begin{aligned}
    \det(\lambda I-A)
    &=(\lambda-p)(\lambda-s)-qr\\
    &=\lambda^2-(p+s)\lambda+(ps-qr)=0.
  \end{aligned}
\]
本書ではこれを\bookterm{matrix}{行列}$A$の\textbf{固有値の方程式}と呼ぶ.\footnote{行列の\bookterm{characteristic}{特性方程式}と呼ぶのが一般的である.本編で「特性方程式」と呼んだ方程式には,この固有値の方程式にあたるものと,不動点の値を求める方程式にあたるものがあるので,付録では呼び分けておく.}
$\lambda$を求めた後,成分の連立方程式から対応する\bookterm{eigenvector}{固有ベクトル}を求められる.
たとえば$A=\begin{pmatrix}2&1\\1&2\end{pmatrix}$では
\[
  \lambda^2-4\lambda+3
  =(\lambda-3)(\lambda-1)=0
\]
だから,\bookterm{eigenvalue}{固有値}は$3,\,1$である.
これは数値例で確かめた倍数と一致する.

\bookterm{recurrence}{漸化式}では,次の項を計算するために記録しておく値の組を\textbf{状態}と呼ぶ.
また,\bookterm{matrix}{行列}式が$0$でないときは,\,$A$を掛けた後の2つの値から掛ける前の2つの値を一意に求め直せる.
このように逆向きに戻せる\bookterm{matrix}{行列}を\textbf{可逆な\bookterm{matrix}{行列}}という.

次節では,本編で引いた定数や作った比・和を,\bookterm{matrix}{行列}の中で見直す.
